\section{Guía de entrada: tres lenguajes para una misma dinámica}
\label{sec:ruta-geometrica-libros}

Este capítulo fija el vocabulario con el que se leerá el reporte. No presupone una
formación previa en topología, pero sí cálculo diferencial e integral, álgebra
lineal y la resolución elemental de ecuaciones diferenciales. La meta no es
memorizar un catálogo de definiciones: es aprender a pasar de una ecuación a
una afirmación geométrica precisa, reconocer sus hipótesis, demostrar lo que
sea accesible y distinguirlo de lo que sólo sugiere una simulación.

El orden bibliográfico rector se presenta, numerado y con criterios de avance,
en la \cref{sec:programa-estudio-bibliografico}. Las actividades de este
capítulo preparan esa ruta y no constituyen un itinerario paralelo.

Los tres títulos adquiridos cumplen funciones diferentes. La ficha
bibliográfica disponible identifica el texto de Amster como una introducción
a métodos de disparo, puntos fijos y grado topológico en problemas no lineales
\cite{Amster2021MetodosTopologicos}. El libro de Nahmad-Achar desarrolla el
lenguaje de topología y geometría diferencial, variedades, formas y Stokes,
con aplicaciones físicas \cite{NahmadAchar2022Topologia}. Los tomos de
Jaurez Rosas, Ortiz Bobadilla, Palma Márquez y Rosales González se orientan
a la teoría geométrica de las ecuaciones diferenciales
\cite{JaurezEtAl2021TGEDI,JaurezEtAl2021TGEDII}. El número
\texttt{8626000002609} es el código comercial del paquete de ambos tomos, no el
ISBN de ninguno de ellos; los ISBN individuales se registran en la
bibliografía.

\begin{remark}[alcance de las referencias de lectura]
No se asignan páginas ni capítulos que no hayan sido verificados directamente
en los ejemplares. La correspondencia siguiente es temática. Si se incorporan
fotografías de los índices, podrá convertirse posteriormente en una tabla
exacta de capítulo, sección y ejercicios sin modificar la cadena matemática.
\end{remark}

\subsection{Tres significados de la palabra topológico}

\begin{definition}[topología del espacio de fases]
Una \emph{topología} sobre un conjunto $X$ es una familia $\mathcal T$ de
subconjuntos, llamados abiertos, que contiene $\varnothing$ y $X$, es cerrada
bajo uniones arbitrarias y bajo intersecciones finitas. La pareja
$(X,\mathcal T)$ es un espacio topológico.
\end{definition}

En dinámica, esta estructura permite decir rigurosamente qué significa estar
cerca, converger, depender continuamente de datos, tener una vecindad de un
equilibrio o separar una cuenca de otra. Una métrica induce una topología, pero
la topología conserva menos información que la distancia: un homeomorfismo
puede deformar una bola en un elipsoide sin romper continuidad, conexidad o
compacidad.

La expresión \emph{métodos topológicos} puede referirse, sin embargo, a tres
carriles distintos:

\begin{enumerate}
 \item \textbf{Geometría diferencial:} variedades, espacios tangentes,
 campos vectoriales, formas diferenciales y métricas. Responde dónde vive el
 sistema y qué estructuras geométricas posee.
 \item \textbf{Dinámica topológica y teoría geométrica de EDO:} flujos,
 equivalencias, conjuntos invariantes, variedades estable e inestable,
 secciones de Poincaré, bifurcaciones y mecanismos de caos. Responde cómo se
 organiza globalmente el conjunto de órbitas.
 \item \textbf{Análisis no lineal topológico:} teoremas de punto fijo, grado
 de Brouwer o Leray--Schauder y continuación. Responde si una solución,
 equilibrio u órbita periódica debe existir aun sin calcularla.
\end{enumerate}

\begin{center}
\begin{tikzpicture}[
 box/.style={draw=Azul,rounded corners,fill=AzulClaro,align=center,
              minimum width=4.4cm,minimum height=1.25cm},
 arr/.style={-{Latex[length=2.5mm]},thick,draw=GrisOscuro},
 node distance=1.7cm and 1.15cm]
 \node[box] (geom) {Variedades y geometría\\campo $f\in T_xM$};
 \node[box,right=of geom] (flow) {Flujo y dinámica\\$\varphi_t:M\to M$};
 \node[box,right=of flow] (top) {Invariantes y topología\\cuencas, retornos, conjugación};
 \node[box,below=of flow] (exist) {Existencia topológica\\punto fijo, grado, continuación};
 \draw[arr] (geom) -- node[above,font=\small]{integra} (flow);
 \draw[arr] (flow) -- node[above,font=\small]{organiza} (top);
 \draw[arr] (exist) -- node[right,font=\small]{garantiza objetos} (flow);
 \draw[arr,dashed] (top) to[bend left=18]
   node[below,font=\small]{formula el problema de ocultedad} (exist);
\end{tikzpicture}

\smallskip
\small\textbf{Diagrama conceptual.} Los tres lenguajes que se enlazan con la
metodología del reporte.
\end{center}

\subsection{Contrato de aprendizaje: definición, hipótesis, prueba y cálculo}

Para cada concepto se seguirá la misma secuencia:

\begin{enumerate}[label=\textbf{G\arabic*.}]
 \item escribir la definición con cuantificadores y espacio ambiente;
 \item listar las hipótesis del teorema antes de emplearlo;
 \item reconstruir el argumento o, cuando exceda el nivel de la guía,
 identificar exactamente qué paso se cita;
 \item estudiar un ejemplo positivo y un contraejemplo;
 \item traducir la afirmación a una operación numérica verificable;
 \item declarar qué conclusión no se obtiene de ese cálculo.
\end{enumerate}

Por ejemplo, una nube aparentemente complicada no es una definición de caos;
un exponente finito positivo no demuestra una herradura; cero contactos en
una muestra no demuestra separación global de cuencas. Esta disciplina conecta
el estudio teórico con los niveles de evidencia L0--L6 del reporte.

\subsection{Diccionario entre el cálculo algebraico y la geometría}

\begin{table}[H]
\centering
\caption{El mismo objeto visto desde cuatro niveles.}
\label{tab:diccionario-algebra-geometria}
\small
\begin{tabularx}{\textwidth}{@{}L{2.7cm}YYY@{}}
\toprule
Objeto & Escritura algebraica & Lectura geométrica & Pregunta topológica \\
\midrule
Campo & $\dot x=f(x)$ & vector $f(x)\in T_xM$ & ¿define un flujo continuo? \\
Equilibrio & $f(e)=0$ & singularidad del campo & ¿qué retrato local conserva un homeomorfismo? \\
Órbita & $\{\varphi_t(x_0)\}$ & curva integral orientada & ¿cuáles son sus conjuntos límite? \\
Linealización & $\dot\xi=J_f(e)\xi$ & acción tangente local & ¿hay conjugación local con el sistema lineal? \\
Ciclo & $\varphi_T(x)=x$ & curva cerrada invariante & ¿es aislado y qué separa en el espacio? \\
Retorno & $P:\Sigma\to\Sigma$ & intersección sucesiva con una sección & ¿qué revela la dinámica iterada de $P$? \\
Cuenca & $\mathcal B(A)$ & conjunto de estados que convergen a $A$ & ¿cómo es su frontera y qué la separa? \\
Ocultedad & $\mathcal B(A)\cap U_e=\varnothing$ & inaccesibilidad desde equilibrios & ¿la separación es local, muestreada o global? \\
\bottomrule
\end{tabularx}
\end{table}

\subsection{Cuatro ejemplos mínimos antes de entrar al caos}

\subsubsection{Un punto silla y sus variedades invariantes}

Considérese
\begin{equation}
 \dot x=x,\qquad \dot y=-y.
 \label{eq:guide-linear-saddle}
\end{equation}
La solución es $\varphi_t(x_0,y_0)=(e^t x_0,e^{-t}y_0)$. Por tanto,
\[
 W^s(0)=\{(0,y):y\in\R\},\qquad
 W^u(0)=\{(x,0):x\in\R\}.
\]
La primera recta contiene exactamente los estados que convergen al origen
cuando $t\to+\infty$; la segunda hace lo mismo cuando $t\to-\infty$. No basta
dibujar flechas: la fórmula demuestra invariancia porque un punto de cada eje
permanece en ese eje para todo tiempo de existencia.

\subsubsection{Una frontera de cuenca calculable}

Para
\begin{equation}
 \dot x=x-x^3=x(1-x^2),
 \label{eq:guide-bistable-line}
\end{equation}
los equilibrios son $-1,0,1$. El signo del campo muestra que
$\mathcal B(-1)=(-\infty,0)$ y $\mathcal B(1)=(0,\infty)$, mientras que
$\{0\}$ es la frontera común. Este ejemplo separa tres conceptos: atractor,
cuenca y frontera. Una imagen del atractor no contiene por sí sola la
información de su cuenca.

\subsubsection{La topología cilíndrica del PLL}

Una fase $\theta$ y $\theta+2\pi$ representan el mismo punto. Por eso el PLL
no vive naturalmente en $\R^2$, sino en $\R\times S^1$. La distancia angular
correcta utiliza
\begin{equation}
 d_{S^1}(\theta_1,\theta_2)
 =\min_{k\in\mathbb Z}|\theta_1-\theta_2+2\pi k|.
 \label{eq:guide-circle-distance}
\end{equation}
Una curva que parece escapar horizontalmente en un rectángulo puede ser una
órbita cerrada sobre el cilindro. Esta observación explica por qué el
expediente PLL emplea retornos módulo $2\pi$, separatriz cilíndrica y una
métrica envuelta.

\subsubsection{Existencia de un ciclo como punto fijo}

Sea $\Sigma$ una sección transversal y $P:\Sigma\to\Sigma$ el primer retorno.
Si $P(z)=z$, la órbita que sale de $z$ vuelve al mismo punto y genera una
órbita periódica. Así, los teoremas de punto fijo del carril de Amster pueden
garantizar ciclos cuando se verifican las hipótesis del operador de retorno.
El argumento no prueba caos: un único punto fijo atractivo corresponde a un
ciclo regular.

\subsection{Mapa de problemas dentro de la ruta única}

La tabla siguiente no cambia el orden oficial: explica qué producto del
proyecto corresponde a cada etapa de la
\cref{sec:programa-estudio-bibliografico}.

\begin{longtable}{@{}L{1.05cm}L{3.0cm}L{4.1cm}L{5.45cm}@{}}
\caption{Secuencia recomendada y producto verificable de cada etapa.}
\label{tab:ruta-estudio-geometrica}\\
\toprule
Etapa & Núcleo conceptual & Fuente temática principal & Producto dentro del proyecto \\
\midrule
\endfirsthead
\toprule
Etapa & Núcleo conceptual & Fuente temática principal & Producto dentro del proyecto \\
\midrule
\endhead
1 & Abiertos, continuidad, compacidad, conexidad y cocientes & Nahmad-Achar & Justificar $S^1=\R/(2\pi\mathbb Z)$ y la métrica cilíndrica del PLL. \\
2 & Variedades, cartas, tangentes y campos & Nahmad-Achar; Teoría geométrica I & Interpretar $\dot x=f(x)$ como sección de $TM$ y reconocer coordenadas físicas. \\
3 & Existencia, unicidad, dependencia continua y flujo & Teoría geométrica I & Identificar dónde Chua PWL sólo es regionalmente $C^1$. \\
4 & Equilibrios, linealización y Hartman--Grobman & Teoría geométrica I--II & Separar espectro, estabilidad local y retrato topológico local. \\
5 & Variedades estable, inestable y central & Teoría geométrica II & Formular las separatrices relevantes de Chua, MAVPD y PLL. \\
6 & Poincaré, monodromía y Floquet & Teoría geométrica II & Reconstruir el multiplicador del ciclo PLL y auditar el retorno MAVPD. \\
7 & Puntos fijos, disparo y grado & Amster & Distinguir existencia de ciclo de su localización numérica. \\
8 & Bifurcaciones y conexiones homoclínicas & Teoría geométrica II & Diseñar continuaciones que registren estabilidad y no sólo amplitud. \\
9 & Caos topológico y dinámica simbólica & Fuentes clásicas complementarias & Comparar herradura o entropía con Lyapunov, FFT y prueba $0$--$1$. \\
10 & Cuencas, fronteras y ocultedad & Este reporte & Sustituir la frase ``se ve oculto'' por un contrato de separación declarado. \\
11 & Semidinámica Caputo en historias & Capítulo fraccionario de esta guía & Diferenciar cuenca funcional, sección de reinicio y continuidad con memoria. \\
\bottomrule
\end{longtable}

\subsection{Actividades guiadas conectadas con los expedientes}

\begin{enumerate}[label=\textbf{A\arabic*.}]
 \item \textbf{Chua PWL.} En cada región, calcule el flujo lineal afín y
 compruebe cómo una trayectoria cruza la superficie $x=\pm1$. Determine qué
 afirmaciones requieren derivada clásica y cuáles sólo continuidad.
 \item \textbf{Kalman--Fitts.} Use el Jacobiano para clasificar el equilibrio y
 describa la cuenca del ciclo candidato. Explique por qué cero contactos en 48
 sondeos no determina su topología global ni prueba ocultedad.
 \item \textbf{MAVPD.} Defina una sección transversal y derive las condiciones
 que evitan cruces tangenciales. Compare el conjunto de cruces interpolados
 del reporte con un verdadero mapa de primer retorno.
 \item \textbf{PLL.} Demuestre que la proyección $\R\to S^1$ identifica fases
 separadas por $2\pi k$. Reconstruya la distancia de
 \cref{eq:guide-circle-distance} y explique por qué una distancia euclídea
 sin envolver puede clasificar mal dos órbitas próximas.
 \item \textbf{Puntos perpetuos.} Verifique que $J_f(p)f(p)=0$ no es invariante
 bajo un cambio de coordenadas no lineal. Discuta qué significaría reemplazar
 la aceleración ordinaria por una aceleración covariante en una variedad.
 \item \textbf{Cuencas.} Para \cref{eq:guide-bistable-line}, muestree una malla
 que no contenga $0$. Explique por qué el experimento aproxima bien ambas
 cuencas, pero no descubre por sí solo la frontera exacta.
 \item \textbf{Lectura crítica.} Tome una figura numérica archivada de uno de
 los expedientes y escriba
 dos columnas: ``lo que muestra'' y ``lo que no demuestra''. Incluya
 acotamiento finito, recurrencia, caos topológico y ocultedad.
\end{enumerate}

\subsection{Autoevaluación conceptual}

Antes de avanzar a la teoría fraccionaria, el lector debería poder responder,
sin recurrir a una figura:

\begin{enumerate}
 \item ¿Cuál es la diferencia entre una trayectoria, una órbita y un flujo?
 \item ¿Qué propiedades conserva un homeomorfismo y cuáles requieren una
 métrica o una estructura diferenciable?
 \item ¿Por qué Hartman--Grobman exige hiperbolicidad?
 \item ¿Por qué un punto fijo de Poincaré representa un ciclo?
 \item ¿Qué diferencia hay entre probar la existencia de una órbita y
 localizar una condición inicial en su cuenca?
 \item ¿Por qué una sección de Poincaré finita no prueba una conjugación
 con un corrimiento simbólico?
 \item ¿Qué parte de la definición de atractor oculto es topológica y qué
 parte del protocolo del reporte es sólo estadística o numérica?
\end{enumerate}

\subsection{Pistas y criterios de autocorrección}

\begin{itemize}
 \item En \textbf{A1}, la continuidad del campo no implica continuidad del
 Jacobiano. La respuesta debe separar las tres regiones y las dos superficies
 de conmutación.
 \item En \textbf{A2}, una clasificación local correcta sólo controla una
 vecindad del equilibrio. La respuesta debe nombrar al menos un objeto global
 que no determine el espectro.
 \item En \textbf{A3}, una sección $h(x)=0$ es transversal cuando
 $\nabla h(x)^{\mathsf T}f(x)\neq0$ en el cruce. Debe explicarse qué falla si
 ese producto se anula.
 \item En \textbf{A4}, el entero $k$ se elige para minimizar la diferencia
 angular. Una verificación correcta prueba también la desigualdad triangular
 o reconoce que se trata de la métrica cociente inducida.
 \item En \textbf{A5}, busque el término Hessiano que aparece en la segunda
 derivada de un cambio no lineal. Omitirlo conduce precisamente a la falsa
 invariancia de los puntos perpetuos.
 \item En \textbf{A6}, el diagrama de signos proporciona una prueba exacta en
 una dimensión. La malla sólo proporciona evidencia finita.
 \item En \textbf{A7}, ``estructura irregular'', ``caos'' y ``ocultedad'' deben
 aparecer en niveles distintos; ninguna figura aislada certifica las dos
 últimas propiedades.
\end{itemize}
